\subsection{Simple Assignment: Creating or Rebinding a Name}

Assignment is the first operation that makes Python variables feel different from C variables. In C, a variable often suggests a typed storage location. In Python, assignment creates or changes a binding between a name and an object.

This continues the model already established earlier: Python variables are not C boxes; they are names that refer to runtime objects. :contentReference[oaicite:0]{index=0}

The basic assignment form is:

\begin{verbatim}
name = object_expression
\end{verbatim}

Python first evaluates the expression on the right side. Then it binds the name on the left side to the resulting object.

\begin{verbatim}
answer = 42

print(answer)
print(type(answer))

print("selected int methods:")
print(f"{'bit_length':<15} {'bit_length' in dir(answer)}")
print(f"{'to_bytes':<15} {'to_bytes' in dir(answer)}")
print(f"{'__add__':<15} {'__add__' in dir(answer)}")

print()
print(f"bit length: {answer.bit_length()}")
print(f"answer + 8: {answer + 8}")
\end{verbatim}

Possible output:

\begin{verbatim}
42
<class 'int'>
selected int methods:
bit_length      True
to_bytes        True
__add__         True

bit length: 6
answer + 8: 50
\end{verbatim}

The name \texttt{answer} does not contain the integer in the C sense. It refers to an integer object. That object has a runtime type, and the type supplies operations such as addition, conversion to string form, and integer-specific methods such as \texttt{bit\_length()}.

\subsubsection{Assignment as Binding, Not Memory Copying}

The assignment:

\begin{verbatim}
answer = 42
\end{verbatim}

does not mean: copy the raw binary value \texttt{42} into a named integer box.

A better conceptual picture is:

\begin{verbatim}
answer  --->  int object 42
\end{verbatim}

The source name \texttt{answer} is a binding label. The object is the runtime value. The object has a type. The name does not have a declared type.

A simple inspection:

\begin{verbatim}
answer = 42

print(f"name:  answer")
print(f"value: {answer}")
print(f"type:  {type(answer)}")
print(f"id:    {id(answer)}")
\end{verbatim}

Possible output:

\begin{verbatim}
name:  answer
value: 42
type:  <class 'int'>
id:    140734432946520
\end{verbatim}

The exact \texttt{id()} value is not important. It is an implementation-level identity number for the object during that run of the program. It lets us observe whether two names refer to the same object.

Assignment can also bind a name to a string object:

\begin{verbatim}
quote = "No soup for you."

print(f"quote: {quote}")
print(f"type:  {type(quote)}")
print(f"id:    {id(quote)}")

print()
print("selected str methods:")
print(f"{'upper':<12} {'upper' in dir(quote)}")
print(f"{'replace':<12} {'replace' in dir(quote)}")
print(f"{'split':<12} {'split' in dir(quote)}")

print()
print(quote.upper())
print(quote.replace("soup", "debugger"))
\end{verbatim}

Possible output:

\begin{verbatim}
quote: No soup for you.
type:  <class 'str'>
id:    2113631854576

selected str methods:
upper        True
replace      True
split        True

NO SOUP FOR YOU.
No debugger for you.
\end{verbatim}

The string object supplies string operations. The name \texttt{quote} only gives us access to that object.

This distinction is essential:

\begin{itemize}
    \item the name is not the object;
    \item the name is not the type;
    \item the name is a binding to an object;
    \item the object carries its runtime type and behavior.
\end{itemize}

\subsubsection{Reassignment: Moving a Name from One Object to Another}

A Python name can be rebound. Rebinding means that the name stops referring to one object and starts referring to another object.

\begin{verbatim}
item = 42

print("first binding")
print(f"value: {item!r}")
print(f"type:  {type(item)}")
print(f"id:    {id(item)}")

print()

item = "It is merely a flesh wound."

print("second binding")
print(f"value: {item!r}")
print(f"type:  {type(item)}")
print(f"id:    {id(item)}")
\end{verbatim}

Possible output:

\begin{verbatim}
first binding
value: 42
type:  <class 'int'>
id:    140734432946520

second binding
value: 'It is merely a flesh wound.'
type:  <class 'str'>
id:    2113631854928
\end{verbatim}

The name \texttt{item} did not transform from an integer box into a string box. The binding changed.

Conceptually:

\begin{verbatim}
first:
    item  --->  int object 42

second:
    item  --->  str object "It is merely a flesh wound."
\end{verbatim}

This is why Python allows this:

\begin{verbatim}
value = 42
value = 3.14
value = "D'oh!"
value = True

print(value)
print(type(value))
\end{verbatim}

Possible output:

\begin{verbatim}
True
<class 'bool'>
\end{verbatim}

Only the final binding remains accessible through the name \texttt{value}. The earlier objects may still exist for a while because of interpreter internals, caching, constants, or other references, but the name \texttt{value} no longer points to them.

A slightly clearer trace:

\begin{verbatim}
thing = 42
print(f"{'step 1':<8} {thing!r:<30} {type(thing)}")

thing = "Chuck Norris can hear sign language."
print(f"{'step 2':<8} {thing!r:<30} {type(thing)}")

thing = 99.5
print(f"{'step 3':<8} {thing!r:<30} {type(thing)}")
\end{verbatim}

Possible output:

\begin{verbatim}
step 1   42                             <class 'int'>
step 2   'Chuck Norris can hear sign language.' <class 'str'>
step 3   99.5                           <class 'float'>
\end{verbatim}

The name stayed the same. The referenced object changed.

This is dynamic binding at the level of assignment. Python does not require the name to keep one declared type for the whole program.

\subsubsection{Object Sharing: Binding Multiple Names to the Same Runtime Object (\texttt{a = b})}

Assignment can also bind a second name to an object that already has a name.

\begin{verbatim}
a = 42
b = a

print(f"a value: {a}")
print(f"b value: {b}")

print()
print(f"a type: {type(a)}")
print(f"b type: {type(b)}")

print()
print(f"a id: {id(a)}")
print(f"b id: {id(b)}")
print(f"a is b: {a is b}")
\end{verbatim}

Possible output:

\begin{verbatim}
a value: 42
b value: 42

a type: <class 'int'>
b type: <class 'int'>

a id: 140734432946520
b id: 140734432946520
a is b: True
\end{verbatim}

The assignment:

\begin{verbatim}
b = a
\end{verbatim}

does not copy the integer object into a new independent object. It binds \texttt{b} to the same object that \texttt{a} already references.

Conceptually:

\begin{verbatim}
a  ----\
       ---> int object 42
b  ----/
\end{verbatim}

The same behavior can be shown with a string:

\begin{verbatim}
line1 = "Nobody expects the Spanish Inquisition!"
line2 = line1

print(f"line1: {line1}")
print(f"line2: {line2}")

print()
print(f"type(line1): {type(line1)}")
print(f"type(line2): {type(line2)}")

print()
print(f"id(line1): {id(line1)}")
print(f"id(line2): {id(line2)}")
print(f"line1 is line2: {line1 is line2}")
\end{verbatim}

Possible output:

\begin{verbatim}
line1: Nobody expects the Spanish Inquisition!
line2: Nobody expects the Spanish Inquisition!

type(line1): <class 'str'>
type(line2): <class 'str'>

id(line1): 2113631855760
id(line2): 2113631855760
line1 is line2: True
\end{verbatim}

Again, there are two names but one object.

If one name is later rebound, the other name is not dragged along with it:

\begin{verbatim}
a = 42
b = a

print("before rebinding")
print(f"a = {a}, id(a) = {id(a)}")
print(f"b = {b}, id(b) = {id(b)}")
print(f"a is b: {a is b}")

print()

a = 100

print("after rebinding a")
print(f"a = {a}, id(a) = {id(a)}")
print(f"b = {b}, id(b) = {id(b)}")
print(f"a is b: {a is b}")
\end{verbatim}

Possible output:

\begin{verbatim}
before rebinding
a = 42, id(a) = 140734432946520
b = 42, id(b) = 140734432946520
a is b: True

after rebinding a
a = 100, id(a) = 140734432948376
b = 42, id(b) = 140734432946520
a is b: False
\end{verbatim}

The line:

\begin{verbatim}
a = 100
\end{verbatim}

does not edit the old integer object \texttt{42}. It rebinds \texttt{a} to a different integer object. The name \texttt{b} still refers to the original object.

Conceptually:

\begin{verbatim}
before:
    a  ----\
           ---> int object 42
    b  ----/

after:
    a  ---> int object 100

    b  ---> int object 42
\end{verbatim}

This example uses integers, which are immutable objects. Immutable means that the object itself cannot be changed in place. Later, when mutable containers are introduced, the distinction between rebinding a name and mutating a shared object will become even more important.

For now, the rule is:

\begin{quote}
Assignment binds a name to an object. Assignment from another name copies the reference, not the object. Reassignment moves one name to another object; it does not move all names that previously shared the old object.
\end{quote}